\section{Related work}
\label{sec:related}

\paragraph{Inverse folding.}
Structured Transformer~\cite{struc_transformer} (2019) introduced the kNN graph
+ autoregressive formulation; GVP-GNN~\cite{gvp} (2021) added geometric vector
perceptrons; ProteinMPNN~\cite{proteinmpnn} (2022) simplified to a
message-passing encoder--decoder with a \emph{random decoding order} and became
the field's default; ESM-IF1~\cite{esm_if1} (2022) showed that 12 M AlphaFold2
structures buy more than architecture (+9.4 pp over GVP-GNN, and +13.3 pp over
the same architecture trained on CATH only in our reading of its tables);
PiFold~\cite{pifold} (2023)
decodes in one shot; LigandMPNN~\cite{ligandmpnn} (2025) conditions on atomic
context. All of these are evaluated on recovery and scTM; none measures
controllability.

\paragraph{Developability-oriented design.}
SolubleMPNN~\cite{soluble_mpnn} (2024) retrains ProteinMPNN on the same PDB
snapshot minus 17,369 transmembrane entries; HyperMPNN~\cite{hypermpnn} (2024)
retrains on thermophile AF2 structures and, in its own measurements, pays a
large expression tax (soluble yield 0.4 vs 25.8 mg/L for the parent model);
MembraneMPNN~\cite{membranempnn} ships conditioned weights
without a paper. Their common structure is a \textbf{prior shift}, which is
irreversible and leaks into every property at once.

\paragraph{Inference-time control.}
Classifier guidance, classifier-free guidance, SMC/Feynman--Kac particle
steering (applied to RFdiffusion in 2025~\cite{fk_rfdiffusion}), and
plug-and-play discrete guidance (ProteinGuide,
2025~\cite{discrete_guidance,proteinguide}) whose AA-ARM update is
\[
  \tilde{p}(x_i{=}a \mid x_S) \;\propto\;
  p_\theta(a \mid x_S)
  \left[
    \frac{p(y \mid x_S \cup a)}{p(y \mid x_S)}
  \right]^{\gamma}.
\]
ZetaDial (2026)~\cite{zetadial} is the closest protein work: a global charge
bias solved per protein, reported to collapse foldability at strong tilt. No
work compares this family against retrained developability models under one
protocol.

\paragraph{Benchmarks.}
ProteinInvBench~\cite{proteininvbench} (NeurIPS 2023 D\&B) standardises
single-chain, multi-chain and de novo tasks; MotifBench~\cite{motifbench} (2025)
standardises motif scaffolding; ProteinGym~\cite{proteingym} and
FLIP~\cite{flip} cover sequence--function. None measures \emph{whether a model
can be steered}, which is the gap this benchmark fills.
